El ser humano se caracteriza por la búsqueda constante de herramientas que mejoren sus condiciones de vida. En los últimos años, estos esfuerzos se han orientado hacia la construcción de máquinas con algoritmos capaces de resolver, de forma automática y rápida, operaciones que resultan tediosas para una persona. Sin embargo, estas máquinas, basadas en algoritmos específicos, presentan una limitación importante: en ocasiones el problema que se busca resolver no requiere un tratamiento algorítmico sino el aprendizaje de patrones a través de la experiencia. 

El ser humano es capaz de resolver múltiples problemas que no pueden expresarse mediante un algoritmo gracias a su experiencia y a su capacidad para memorizar y asociar hechos. Así, parece lógico que una forma de aproximarse a este tipo de problemas consista en construir sistemas capaces de reproducir dicha característica humana. En este sentido, las redes neuronales constituyen un intento simplificado de simular el funcionamiento del cerebro humano en el tratamiento de la información, cuya unidad básica de procesamiento se inspira en la célula fundamental del sistema nervioso: la neurona.

Una característica esencial del cerebro humano es su capacidad de aprendizaje, entendida como la habilidad de resolver problemas que inicialmente no podían abordarse, mediante la adquisición de nueva información sobre ellos \citep{matich2001redes, salas2004redes, tablada2009redes, izaurieta2000redes}.

Por lo tanto, las redes neuronales se componen de unidades de procesamiento que intercambian datos o información. Se utilizan para reconocer patrones en conjuntos de datos, como imágenes, manuscritos o secuencias temporales, con el objetivo de aprender de ellos y mejorar su desempeño \citep{matich2001redes}.

Debido a sus fundamentos, las redes neuronales artificiales presentan un gran número de características semejantes a las del cerebro, y por ello una gran cantidad de ventajas por sobre los programas basados en algoritmos \cite{McCullochPitts1943}. Por ejemplo, son capaces de aprender a través de la experiencia, de generalizar de una gran serie de casos específicos a nuevos casos, de abstraer características esenciales de bancos de información, desechando la información irrelevante, etc. 

